% Reproducible 2-D random-walk Kalman example from the chapter text.
\begin{tikzpicture}
  \begin{axis}[
    width=.92\linewidth,height=82mm,
    xmin=-.45,xmax=8.45,ymin=-.65,ymax=5.25,
    xlabel={横向位置},ylabel={纵向位置},
    xtick={0,2,4,6,8},ytick={0,1,2,3,4,5},
    tick label style={font=\figurefont\footnotesize},
    label style={font=\figurefont\footnotesize},
    axis line style={FigureMuted},tick style={FigureMuted},
    grid=major,grid style={FigureGrid,line width=.5pt},
    legend style={font=\figurefont\footnotesize,draw=FigureGrid,
      fill=none,at={(.02,.98)},anchor=north west,
      cells={anchor=west}},
    axis equal image,clip=false]
    \addplot[FigureInk,densely dashed,line width=1.1pt,mark=*,mark size=1.4pt]
      coordinates {
        (0,0) (1,.5) (2,1.1) (3,1.4) (4,2.2)
        (5,2.6) (6,3.4) (7,3.7) (8,4.4)
      };
    \addlegendentry{真值}
    \addplot[FigureMuted,line width=.8pt,mark=x,mark size=2.2pt]
      coordinates {
        (.2,-.3) (1.4,.2) (1.6,1.5) (3.5,1.0) (3.7,2.7)
        (5.4,2.2) (5.7,3.9) (7.5,3.3) (7.8,4.8)
      };
    \addlegendentry{观测}
    \addplot[FigureBlue,line width=1pt,mark=o,mark size=1.7pt]
      coordinates {
        (0,0) (.094,-.171) (.998,.085) (1.427,1.079)
        (2.906,1.023) (3.473,2.202) (4.848,2.201)
        (5.456,3.396) (6.915,3.328)
      };
    \addlegendentry{一步预测}
    \addplot[FigureRed,line width=1.25pt,mark=square*,mark size=1.5pt]
      coordinates {
        (.094,-.171) (.998,.085) (1.427,1.079) (2.906,1.023)
        (3.473,2.202) (4.848,2.201) (5.456,3.396)
        (6.915,3.328) (7.546,4.363)
      };
    \addlegendentry{滤波}
    \addplot[FigureYellow,line width=1.35pt,mark=triangle*,mark size=1.6pt]
      coordinates {
        (.338,-.024) (1.261,.405) (1.935,1.176) (3.206,1.406)
        (3.955,2.313) (5.157,2.576) (5.926,3.467)
        (7.096,3.636) (7.546,4.363)
      };
    \addlegendentry{RTS平滑}
    \addplot[FigureRed,fill=FigureRedFill,fill opacity=.22,
      line width=.8pt,domain=0:360,samples=121]
      ({1.427+.566*cos(x)},{1.079+.459*sin(x)});
    \addplot[FigureRed,fill=FigureRedFill,fill opacity=.22,
      line width=.8pt,domain=0:360,samples=121]
      ({4.848+.567*cos(x)},{2.201+.459*sin(x)});
    \addplot[FigureRed,fill=FigureRedFill,fill opacity=.22,
      line width=.8pt,domain=0:360,samples=121]
      ({7.546+.567*cos(x)},{4.363+.459*sin(x)});
    \node[anchor=south east,text=FigureInk,fill=none,inner sep=1pt]
      at (axis cs:4.73,2.57) {$1\sigma$滤波椭圆};
    \node[font=\figurefont\scriptsize,anchor=south west,text=FigureInk]
      at (axis cs:.094,-.171) {$t=0$};
    \node[font=\figurefont\scriptsize,anchor=south west,text=FigureInk]
      at (axis cs:3.473,2.202) {$t=4$};
    \node[font=\figurefont\scriptsize,anchor=south west,text=FigureInk]
      at (axis cs:7.546,4.363) {$t=8$};
  \end{axis}
\end{tikzpicture}
